\section{Related Work}

Generative models have significantly advanced our ability to synthesize realistic medical images, helping to alleviate data‑scarcity and class‑imbalance issues and improving the performance of downstream tasks such as segmentation and diagnosis. Existing research has explored GAN-based methods for generating 3D medical images. \citet{kwon2019generation} combine an auto-encoder with a GAN model to generate 3D brain MRI samples, while \citet{hong20213d} tackle this problem by extending StyleGAN2~\citep{karras2020analyzing}. %in their 3D‑StyleGAN work. 
HA-GAN~\citep{sun2022hierarchical} instead employs a hierarchical patch‑based generator and discriminator to produce high‑resolution 3D medical images. 

Despite the promising works, GAN-based models still suffer from a range of issues including mode collapse, substantial memory requirements and low generation quality. Consequently, recent research has increasingly focused on diffusion models as a promising alternative, due to their increased training stability and higher fidelity in generated samples. 3D-DDPM~\citep{dorjsembe2022three} introduce denoising diffusion probabilistic models tailored for 3D medical data to generate high-resolution brain MRI scans. However, building diffusion models directly on 3D medical data incurs significant computational and memory costs, limiting scalability to high-resolution volumes. Following this, several new models propose corresponding improvements. \citet{khader2023medical} combine a VQ-GAN-based latent space with a denoising diffusion probabilistic model to achieve high-quality synthesis of multi-modal 3D medical images, effectively enhancing the clarity and diversity of the generated images. Wavelet Diffusion Model (WDM)~\citep{friedrich2024wdm} improves efficiency by applying diffusion to wavelet-decomposed images, enabling high-resolution synthesis with reduced memory demands. TCAM-Diff~\citep{zhang2025tcam} incorporates a triplane-aware cross-attention mechanism that efficiently generates high-resolution 3D medical images while significantly reducing memory usage. However, all these models are based on variants of the convolutional U-Net architecture~\citep{ronneberger2015u}, and excel at capturing local features but are limited in modeling crucial long-range dependencies. 

To address this limitation \rebuttal{in 2D tasks}, Diffusion Transformer (DiT)~\citep{Peebles2022DiT} models have recently been proposed, replacing convolutional patch embedders with Transformer blocks to better capture global context and long-range dependencies in image generation. \rebuttal{DiT has demonstrated state-of-the-art performance in 2D image synthesis and has also shown promise when adapted to sparse 3D settings such as voxelized point clouds~\citep{mo2023ditd}. However, despite these advances, we found that DiT still encounters substantial challenges in dense 3D scenarios, including unstable training behavior and optimization difficulties, in addition to the high computational cost required for scaling.} Motivated by this, our work introduces a Residual Diffusion Transformer that decomposes the generation process into two stages: a lightweight MLP-based local denoiser operating on 3D patches, followed by a Transformer module that models and refines the global residuals. This two-stage design effectively balances computational efficiency with enhanced capability to capture both fine local details and global structural information, enabling scalable and high-fidelity synthesis of high-resolution 3D medical volumes.